\section{مبانی نظری}

\subsection{بیماری دیستروفی عضلانی دوشن}

دیستروفی عضلانی دوشن (\lr{Duchenne Muscular Dystrophy - DMD}) یکی از شایع‌ترین و شدیدترین بیماری‌های ژنتیکی نوروموسکولار در کودکان است که عمدتاً پسران را درگیر می‌کند و به‌صورت وابسته به کروموزوم X به ارث می‌رسد \cite{Emery2002DMD}. این بیماری معمولاً ناشی از جهش در ژن دیستروفین است که می‌تواند منجر به فقدان یا کاهش شدید پروتئین دیستروفین در فیبرهای عضلانی شود \cite{Hoffman1987Dystrophin}. دیستروفین نقش مهمی در پایداری مکانیکی غشای سلول‌های عضلانی دارد و اختلال در آن موجب آسیب‌پذیری فیبرهای عضلانی در برابر تنش‌های مکانیکی ناشی از انقباض می‌گردد.

از منظر بالینی، نشانه‌های بیماری غالباً در سال‌های ابتدایی کودکی ظاهر می‌شوند و می‌توانند شامل تأخیر در مهارت‌های حرکتی، دشواری در دویدن و بالا رفتن از پله‌ها، الگوی راه رفتن غیرطبیعی و زمین خوردن‌های مکرر باشند \cite{Bushby2010Diagnosis}. یکی از نشانه‌های شناخته‌شده در بسیاری از بیماران، علامت 
\lr{Gowers' sign}
\footnote{الگوی جبرانی برخاستن از زمین که در آن بیمار با کمک دست‌ها و بالا کشیدن تدریجی تنه از روی ران‌ها به وضعیت ایستاده می‌رسد.} 
است که به ضعف عضلات کمربند لگنی نسبت داده می‌شود.

با پیشرفت بیماری، افت عملکرد حرکتی به‌صورت تدریجی رخ می‌دهد و در بسیاری از بیماران در سنین نوجوانی از دست رفتن توانایی راه رفتن مستقل مشاهده می‌شود. همچنین در مراحل پیشرفته‌تر، درگیری سیستم تنفسی و قلبی می‌تواند بر کیفیت زندگی و پیامدهای بالینی اثرگذار باشد \cite{Eagle2002Survival}. ازاین‌رو، ارزیابی و پایش دقیق شاخص‌های عملکرد موتوری برای بررسی روند بیماری، ارزیابی اثربخشی درمان و برنامه‌ریزی توان‌بخشی ضروری است \cite{Bushby2010Diagnosis}.

\subsection{مقیاس NSAA و اصول امتیازدهی}

مقیاس \lr{North Star Ambulatory Assessment (NSAA)} یکی از ابزارهای استاندارد ارزیابی عملکرد حرکتی در بیماران سرپایی مبتلا به دیستروفی عضلانی دوشن است که به‌طور گسترده در کارآزمایی‌های بالینی و مطالعات پایش طولی مورد استفاده قرار می‌گیرد \cite{Scott2009NSAA}. این مقیاس به‌طور اختصاصی برای بیماران DMD طراحی شده و هدف آن اندازه‌گیری توانایی انجام فعالیت‌های عملکردی روزمره مرتبط با راه رفتن و تعادل است.

NSAA شامل ۱۷ فعالیت عملکردی است که طیفی از حرکات پایه تا فعالیت‌های دینامیک‌تر مانند پریدن و دویدن را دربرمی‌گیرد. هر فعالیت بر اساس یک مقیاس ترتیبی سه‌سطحی امتیازدهی می‌شود:

\begin{itemize}
	\item نمره 2: اجرای طبیعی و مستقل حرکت بدون استفاده از الگوهای جبرانی؛
	\item نمره 1: اجرای مستقل همراه با جبران حرکتی یا تغییر در الگوی طبیعی حرکت؛
	\item نمره 0: ناتوانی در انجام مستقل فعالیت.
\end{itemize}

حداکثر امتیاز قابل کسب در این مقیاس ۳۴ است که نشان‌دهنده عملکرد کامل حرکتی است. از آنجا که این مقیاس ماهیتی ترتیبی دارد، فاصله بین نمرات لزوماً بیانگر اختلاف عددی خطی در کیفیت حرکت نیست، بلکه نشان‌دهنده طبقه‌بندی سطوح عملکردی است.

مطالعات ارزیابی پایایی و روایی این ابزار نشان داده‌اند که NSAA از اعتبار بالایی برای پایش تغییرات عملکردی بیماران DMD برخوردار است، با این حال وابستگی آن به مشاهده مستقیم و قضاوت ارزیاب می‌تواند منجر به تغییرپذیری بین‌ارزیاب شود \cite{Mayhew2013Reliability}. این موضوع به‌ویژه در مواردی که تمایز بین اجرای «طبیعی» و اجرای «اصلاح‌شده» ظریف است، اهمیت بیشتری پیدا می‌کند.

از منظر مهندسی تحلیل حرکت، بسیاری از معیارهای کیفی NSAA—مانند استفاده از حرکات جبرانی، تغییر در دامنه حرکتی یا کاهش سرعت اجرای حرکت—می‌توانند با تغییر در پارامترهای سینماتیکی قابل‌اندازه‌گیری همراه باشند. بنابراین، کمی‌سازی این شاخص‌ها با استفاده از حسگرهای پوشیدنی می‌تواند به توسعه یک چارچوب داده‌محور مکمل برای فرآیند نمره‌دهی کمک نماید.

\subsection{حسگرهای اینرسیایی و کاربرد آن‌ها در تحلیل حرکت}

حسگرهای اینرسیایی که معمولاً در قالب واحدهای \lr{Inertial Measurement Unit (IMU)} عرضه می‌شوند، ابزارهایی پوشیدنی برای اندازه‌گیری پارامترهای سینماتیکی حرکت هستند. این واحدها عموماً شامل شتاب‌سنج و ژیروسکوپ سه‌محوره بوده و امکان ثبت شتاب خطی و سرعت زاویه‌ای را در راستای سه محور مختصاتی فراهم می‌کنند \cite{Sabatini2011IMU}.

شتاب‌سنج (\lr{Accelerometer}) میزان شتاب خطی وارد بر حسگر را اندازه‌گیری می‌کند که شامل مؤلفه ناشی از حرکت و مؤلفه گرانشی است. بنابراین، در تحلیل حرکت لازم است اثر شتاب گرانشی از داده‌ها تفکیک یا جبران شود تا مؤلفه‌های دینامیکی حرکت استخراج گردند \cite{Madgwick2011SensorFusion}. ژیروسکوپ (\lr{Gyroscope}) نیز سرعت زاویه‌ای حول محورهای مختصاتی را ثبت می‌کند و امکان تحلیل چرخش اندام‌ها و تغییرات وضعیت زاویه‌ای را فراهم می‌سازد.

در سال‌های اخیر، استفاده از IMUها در تحلیل حرکت انسان به‌طور گسترده‌ای افزایش یافته است، زیرا این حسگرها نسبت به سیستم‌های موشن‌کپچر اپتیکال دارای مزایایی نظیر هزینه کمتر، قابلیت حمل، عدم وابستگی به شرایط نوری و امکان استفاده در محیط‌های غیرآزمایشگاهی هستند \cite{MuroDeLaHerran2014Wearable}. این ویژگی‌ها آن‌ها را به گزینه‌ای مناسب برای کاربردهای بالینی و پایش بیماران در محیط واقعی تبدیل کرده است.

با این حال، داده‌های حاصل از IMUها ممکن است تحت تأثیر نویز اندازه‌گیری، خطای کالیبراسیون و رانش انتگرالی\footnote{افزایش تدریجی خطا در تخمین زاویه که ناشی از انتگرال‌گیری طولانی‌مدت از سرعت زاویه‌ای است.} قرار گیرند. این موضوع به‌ویژه در کاربردهایی که نیاز به تخمین وضعیت زاویه‌ای دقیق دارند، اهمیت بیشتری پیدا می‌کند. ازاین‌رو، استفاده از روش‌های مناسب پیش‌پردازش و فیلترگذاری برای افزایش دقت و پایداری داده‌ها ضروری است.

\subsection{پردازش سیگنال در تحلیل حرکت}

داده‌های خام حاصل از حسگرهای اینرسیایی معمولاً شامل نویز اندازه‌گیری، مؤلفه‌های ناخواسته و تغییرات غیرمرتبط با فعالیت هدف هستند. بنابراین پیش از انجام هرگونه تحلیل کمی، لازم است فرآیند پیش‌پردازش سیگنال انجام شود. یکی از منابع اصلی در زمینه تحلیل بیومکانیکی حرکت، بر اهمیت فیلترگذاری مناسب برای حذف نویز و افزایش نسبت سیگنال به نویز تأکید دارد \cite{Winter2009Biomechanics}.

یکی از روش‌های رایج در تحلیل حرکت، استفاده از فیلترهای دیجیتال پایین‌گذر مانند فیلتر باترورث (\lr{Butterworth filter}) است. این فیلتر به دلیل پاسخ فرکانسی یکنواخت در باند عبور و عدم ایجاد نوسانات شدید در سیگنال، در کاربردهای بیومکانیکی متداول است \cite{Winter2009Biomechanics}. انتخاب فرکانس قطع مناسب باید با توجه به ماهیت حرکت و نرخ نمونه‌برداری انجام شود.

علاوه بر فیلترگذاری، قطعه‌بندی زمانی (\lr{Segmentation}) نقش مهمی در تحلیل فعالیت‌های حرکتی دارد. در این فرآیند، سیگنال پیوسته به بازه‌هایی متناظر با هر فعالیت مشخص تقسیم می‌شود تا امکان استخراج ویژگی‌های مرتبط با هر حرکت فراهم گردد. همچنین نرمال‌سازی زمانی\footnote{فرآیندی که طی آن طول زمانی اجرای حرکات مختلف به یک مقیاس مشترک تبدیل می‌شود تا امکان مقایسه بین افراد با سرعت‌های اجرای متفاوت فراهم گردد.} به مقایسه بین آزمودنی‌ها یا اجرای‌های مختلف یک فعالیت کمک می‌کند.

پس از پیش‌پردازش، مرحله استخراج ویژگی (\lr{Feature Extraction}) انجام می‌شود. ویژگی‌های حوزه زمان مانند میانگین، انحراف معیار، دامنه تغییرات و زمان اجرای حرکت، و ویژگی‌های حوزه فرکانس مانند توان طیفی یا مؤلفه‌های غالب فرکانسی، از جمله شاخص‌هایی هستند که در مطالعات تحلیل فعالیت انسانی مورد استفاده قرار گرفته‌اند \cite{Bulling2014HumanActivity}. انتخاب ویژگی‌های مناسب نقش تعیین‌کننده‌ای در عملکرد الگوریتم‌های یادگیری ماشین دارد و باید با توجه به ماهیت مسئله و نوع برچسب‌های مورد پیش‌بینی انجام شود.

\subsection{کاربرد یادگیری ماشین در تحلیل حرکت}

یادگیری ماشین (\lr{Machine Learning}) مجموعه‌ای از روش‌های محاسباتی است که به سیستم‌ها امکان می‌دهد الگوهای موجود در داده‌ها را یاد گرفته و بر اساس آن پیش‌بینی یا طبقه‌بندی انجام دهند \cite{Bishop2006PatternRecognition}. در تحلیل حرکت انسان، این روش‌ها به‌طور گسترده برای شناسایی فعالیت‌ها، طبقه‌بندی الگوهای حرکتی و پیش‌بینی شاخص‌های عملکردی به کار گرفته شده‌اند.

در چارچوب یادگیری نظارت‌شده (\lr{Supervised Learning})، مدل با استفاده از داده‌های برچسب‌دار آموزش می‌بیند؛ به این معنا که برای هر نمونه داده، خروجی صحیح (مانند نوع فعالیت یا نمره عملکردی) مشخص است. پس از آموزش، مدل می‌تواند برای داده‌های جدید، خروجی متناظر را پیش‌بینی کند. روش‌هایی مانند ماشین بردار پشتیبان (\lr{Support Vector Machine})، درخت تصمیم و شبکه‌های عصبی مصنوعی از جمله الگوریتم‌هایی هستند که در مطالعات تحلیل فعالیت انسانی کاربرد گسترده‌ای داشته‌اند \cite{Lara2013Survey}.

در زمینه تحلیل داده‌های حسگرهای پوشیدنی، یادگیری ماشین به‌عنوان ابزاری برای استخراج الگوهای پیچیده و غیرخطی از داده‌های چندبعدی شناخته می‌شود. ویژگی‌های استخراج‌شده از سیگنال‌های شتاب و سرعت زاویه‌ای می‌توانند به‌عنوان ورودی مدل قرار گرفته و رابطه آن‌ها با شاخص‌های عملکردی مدل‌سازی شود.

در پژوهش حاضر، مسئله پیش‌بینی نمره حرکات بر اساس مقیاس NSAA را می‌توان به‌صورت یک مسئله طبقه‌بندی یا رگرسیون
\footnote{
	در حالت طبقه‌بندی، مدل سطوح گسسته نمره (۰، ۱، ۲) را پیش‌بینی می‌کند؛ در حالی‌که در رگرسیون، خروجی مدل به‌صورت یک مقدار عددی پیوسته تخمین زده می‌شود.
}
	 در نظر گرفت. انتخاب چارچوب مناسب به ماهیت داده‌ها و نحوه تعریف مسئله وابسته است.

به‌کارگیری رویکرد یادگیری ماشین در این زمینه می‌تواند به کاهش وابستگی به قضاوت ذهنی ارزیاب، افزایش تکرارپذیری و ارائه یک چارچوب کمی برای نمره‌دهی حرکات کمک نماید.
